% Источники: 07_Программные_требования/DOC-06.6_Конфигурации_и_модерация.md, 07_Программные_требования/Конфигурации_и_модерация.md
% Генератор: docs/build_pdf.py
% Шаблон печатного представления. Содержимое подставляет docs/build_pdf.py.
\documentclass[12pt,a4paper]{article}
\usepackage{iftex}
\ifPDFTeX
  \usepackage[utf8]{inputenc}
  \usepackage{cmap}
  \usepackage[T2A]{fontenc}
  % Tempora retains Times-like text and provides Cyrillic bold/italic shapes.
  \usepackage{tempora,newtxmath}
\else
  \usepackage{fontspec}
  \IfFontExistsTF{Times New Roman}
    {\setmainfont{Times New Roman}}
    {\setmainfont{Libertinus Serif}}
  \IfFontExistsTF{Arial}{\setsansfont{Arial}}{\setsansfont{Latin Modern Sans}}
  \IfFontExistsTF{Menlo}{\setmonofont{Menlo}[Scale=MatchLowercase]}
    {\setmonofont{Latin Modern Mono}[Scale=MatchLowercase]}
\fi
\usepackage[russian]{babel}
\usepackage{geometry}
\geometry{left=25mm,right=15mm,top=20mm,bottom=20mm}
\usepackage{setspace}
\onehalfspacing
\usepackage{indentfirst}
\setlength{\parindent}{1.25cm}
\setlength{\parskip}{2pt}
\setlength{\emergencystretch}{3em}
\usepackage{graphicx,booktabs,longtable,array,tabularx,calc,etoolbox}
\newcounter{none} % Pandoc 3.9 uses this counter for unnumbered longtables.
\usepackage{enumitem}
\setlist[itemize]{leftmargin=1.25cm,itemsep=2pt,topsep=4pt}
\setlist[enumerate]{leftmargin=1.25cm,itemsep=2pt,topsep=4pt}
\newcommand{\tightlist}{\setlength{\itemsep}{0pt}\setlength{\parskip}{0pt}}
\AtBeginEnvironment{longtable}{\small\singlespacing}
\setlength{\LTpre}{8pt}
\setlength{\LTpost}{8pt}
\setlength{\tabcolsep}{5pt}
\setlength{\LTleft}{0pt}
\setlength{\LTright}{0pt}
\usepackage{xurl}
\usepackage[unicode,colorlinks=true,linkcolor=black,urlcolor=blue]{hyperref}
\usepackage{bookmark}
\hypersetup{pdftitle={Шаблоны наборов задач и модерация},pdfauthor={Самбьюи А.Ж.},pdfsubject={Криптотренажёр mephi-12}}
\urlstyle{same}
\pagestyle{plain}
\setcounter{secnumdepth}{3}
\setcounter{tocdepth}{2}
\listfiles

\begin{document}
\hypersetup{pageanchor=false}
\begin{titlepage}
  \thispagestyle{empty}
  \centering
  {\large НИЯУ МИФИ\par}
  \vspace*{3.5cm}
  {\LARGE\bfseries Шаблоны наборов задач и модерация\par}
  \vspace{0.5cm}
  {\large DOC-06.6\par}
  \vspace{0.6cm}
  {\large Платформа для изучения и практики курса по криптографии\par}
  {\large Криптотренажёр mephi-12\par}
  \vspace{0.7cm}
  {\large Версия 1.0\par}
  {\large Редакция от 25.11.2026\par}
  \vspace{1cm}
  \begin{tabularx}{\textwidth}{@{}p{3.4cm}Xp{4cm}@{}}
    \toprule
    Роль & ФИО & Подпись / Дата \\
    \midrule
    Разработал & Самбьюи А.Ж. & \\
    \addlinespace[6pt]
    Согласовано & & \\
    \addlinespace[6pt]
    Утверждено & & \\
    \bottomrule
  \end{tabularx}
  \vfill
  {\large Москва, 2026\par}
\end{titlepage}
\hypersetup{pageanchor=true}

\section*{Лист регистрации изменений}
\addcontentsline{toc}{section}{Лист регистрации изменений}
\begingroup\small\singlespacing
\noindent
\begin{tabularx}{\textwidth}{@{}p{1.2cm}p{2cm}p{2.5cm}Xp{1.8cm}@{}}
\toprule
Версия & Дата & Автор & Содержание изменения & Утвердил \\
\midrule
1.0 & 25.11.2026 & Самбьюи А.Ж. & Подготовлен кандидат v1.0; сверены модель, API, кейсы и трассировка, уточнены права подготовки банка и черновика. Принятие ожидается. & \\
\bottomrule
\end{tabularx}
\endgroup
\clearpage
\begingroup\small\singlespacing
\tableofcontents
\endgroup
\clearpage

\hypertarget{doc-066-шаблоны-наборов-задач-и-модерация}{}

{\def\LTcaptype{none} % do not increment counter
\begin{longtable}[]{@{}
  >{\raggedright\arraybackslash}p{(\linewidth - 2\tabcolsep) * \real{0.3200}}
  >{\raggedright\arraybackslash}p{(\linewidth - 2\tabcolsep) * \real{0.6800}}@{}}
\toprule\noalign{}
\begin{minipage}[b]{\linewidth}\raggedright
Реквизит
\end{minipage} & \begin{minipage}[b]{\linewidth}\raggedright
Значение
\end{minipage} \\
\midrule\noalign{}
\endhead
\bottomrule\noalign{}
\endlastfoot
Система & Платформа для изучения и практики курса по криптографии в МИФИ \\
Редакция & 1.0 от 25.11.2026 \\
Владелец подсистемы & Самбьюи А.Ж. \\
Исходная редакторская подготовка требований & Анохин Г.А., сводная редакция DOC-06 от 20.11.2026 \\
Основание & Задачи 2.11 и 3.2; SR-CFG-001--SR-CFG-008 действующего DOC-05 \\
Статус & Кандидат v1.0 для рассмотрения; инспекция и принятие не подтверждены \\
\end{longtable}
}

\section{Назначение и границы}\label{1-ux43dux430ux437ux43dux430ux447ux435ux43dux438ux435-ux438-ux433ux440ux430ux43dux438ux446ux44b}

\hypertarget{1-назначение-и-границы}{}

Подсистема предоставляет шаблоны самостоятельной практики, подготовку контрольных мероприятий, публикацию в расписании, авторские задания с вариантами и просмотр учебного прогресса. Работа относится к домену CFG. Персональная сессия и обработка ответа принадлежат домену SES; подтверждение личности и назначение преподавателю групп --- домену AUTH.

В состав печатного DOC-06.6 входит {\hyperlink{конфигурации-контрольные-мероприятия-и-модерация}{общая глава с определениями SWR-CFG}}. Она используется также в сводном DOC-06, поэтому редактирование одной обязанности выполняется в одном Markdown-источнике. Идентификаторы и предметные решения сохранены. {\href{DOC-06\_Программные\_требования.pdf\#приложение-б-физическая-модель-данных}{Физическая модель данных}} подготовлена как приложение сводного DOC-06; DOC-06.6 использует её по ссылке.

\section{Данные на границах подсистемы}\label{2-ux434ux430ux43dux43dux44bux435-ux43dux430-ux433ux440ux430ux43dux438ux446ux430ux445-ux43fux43eux434ux441ux438ux441ux442ux435ux43cux44b}

\hypertarget{2-данные-на-границах-подсистемы}{}

{\def\LTcaptype{none} % do not increment counter
\begin{longtable}[]{@{}
  >{\raggedright\arraybackslash}p{(\linewidth - 2\tabcolsep) * \real{0.3200}}
  >{\raggedright\arraybackslash}p{(\linewidth - 2\tabcolsep) * \real{0.6800}}@{}}
\toprule\noalign{}
\begin{minipage}[b]{\linewidth}\raggedright
Объект
\end{minipage} & \begin{minipage}[b]{\linewidth}\raggedright
Сведения, необходимые подсистеме
\end{minipage} \\
\midrule\noalign{}
\endhead
\bottomrule\noalign{}
\endlastfoot
Шаблон и его назначение группе & Имя, упорядоченные позиции типов и признак \texttt{enabled} отдельно для пары «группа --- шаблон». Повторение типа означает отдельную позицию \\
Черновик мероприятия & UUID \texttt{assessment\allowbreak{}Id}, автор, название, состояние \texttt{DRAFT}, позиции и необязательный до публикации выбор \texttt{repeat\allowbreak{}Allowed} \\
Опубликованное мероприятие & Назначенная группа, \texttt{opensAt}, \texttt{closesAt}, явное правило повтора, закреплённый порядок позиций и редакции заданий. Состояние --- \texttt{PUBLISHED} \\
Авторское задание & Принадлежащие преподавателю варианты; в каждой паре условие \texttt{statement} и эталон \texttt{expected\allowbreak{}Answer}. Студент получает только открытое условие \\
Выдача личного набора & Связь с мероприятием либо самостоятельной практикой, позиции и закреплённые варианты. UUID, сроки и состояния сессии формирует подсистема сессий \\
Представление прогресса & Владелец и идентификатор прохождения, сохранённые состояния и предметные результаты. Право чтения не даёт права отправлять ответы за студента \\
\end{longtable}
}

Табличное описание задаёт логические обязанности, а не новый SQL-контракт. Общие требования хранения и редакторское описание модели находятся в {\href{DOC-06\_Программные\_требования.pdf\#приложение-а-модель-данных}{DOC-06}}.

\section{Интерфейсы и ответственность}\label{3-ux438ux43dux442ux435ux440ux444ux435ux439ux441ux44b-ux438-ux43eux442ux432ux435ux442ux441ux442ux432ux435ux43dux43dux43eux441ux442ux44c}

\hypertarget{3-интерфейсы-и-ответственность}{}

Подтверждённая личность и текущие назначения групп поступают из подсистемы учётных записей. Для собственного банка и черновика проверяются роль преподавателя и принадлежность объекта; отсутствие назначенной группы не запрещает их подготовку. При публикации и чтении результатов дополнительно требуется действующее назначение конкретной группы. Общий административный просмотр не предоставляется роли преподавателя автоматически.

{\def\LTcaptype{none} % do not increment counter
\begin{longtable}[]{@{}
  >{\raggedright\arraybackslash}p{(\linewidth - 4\tabcolsep) * \real{0.2028}}
  >{\raggedright\arraybackslash}p{(\linewidth - 4\tabcolsep) * \real{0.3853}}
  >{\raggedright\arraybackslash}p{(\linewidth - 4\tabcolsep) * \real{0.4119}}@{}}
\toprule\noalign{}
\begin{minipage}[b]{\linewidth}\raggedright
Граница
\end{minipage} & \begin{minipage}[b]{\linewidth}\raggedright
Обязанность CFG
\end{minipage} & \begin{minipage}[b]{\linewidth}\raggedright
Смежная обязанность
\end{minipage} \\
\midrule\noalign{}
\endhead
\bottomrule\noalign{}
\endlastfoot
Каталог шаблонов & Отобрать шаблоны подтверждённой группы, сохранить порядок и признаки доступности & AUTH подтверждает группу пользователя; SES проверяет допустимость создания личного набора \\
Подготовка и публикация & Сохранить собственный черновик, проверить состав и назначение, зафиксировать опубликованную редакцию & AUTH предоставляет актуальные назначения групп; хранение обеспечивает согласованную запись \\
Авторское задание & Сохранить пары вариантов, предоставить редакцию, закреплённую для публикации & SES связывает выданный вариант с задачей и применяет общие правила попытки \\
Старт по мероприятию & Определить доступность опубликованной работы и выбор повтора & SES обеспечивает единственную активную сессию, UUID, срок и запись созданного набора \\
Просмотр & Применить правила административного или преподавательского допуска & SES определяет представление состояния и сохранённого результата по общему контракту чтения \\
\end{longtable}
}

Пути, поля запросов и HTTP-ответы определены в {\href{../09\_Архитектура\_и\_API/DOC-08\_Справочник\_API.pdf\#doc-08-справочник-api-платформы}{DOC-08}}. Правила сроков, попыток и восстановления определены в {\href{DOC-06.2\_Сессии\_задач\_и\_планировщик.pdf\#doc-062-сессии-задач-и-планировщик}{DOC-06.2}} и не переопределяются этим документом. Время новой личной сессии не отсчитывается заново при чтении, повторной доставке ответа или восстановлении.

\section{Трассировка и условия проверки}\label{4-ux442ux440ux430ux441ux441ux438ux440ux43eux432ux43aux430-ux438-ux443ux441ux43bux43eux432ux438ux44f-ux43fux440ux43eux432ux435ux440ux43aux438}

\hypertarget{4-трассировка-и-условия-проверки}{}

Все требования этой подсистемы обязательны. Источники и наблюдаемые результаты указаны у определений SWR-CFG. Следующие ссылки относятся к реально подготовленным SR и кейсам DOC-10, а не к результатам исполнения тестов.

{\def\LTcaptype{none} % do not increment counter
\begin{longtable}[]{@{}
  >{\raggedright\arraybackslash}p{(\linewidth - 4\tabcolsep) * \real{0.3333}}
  >{\raggedright\arraybackslash}p{(\linewidth - 4\tabcolsep) * \real{0.3333}}
  >{\raggedright\arraybackslash}p{(\linewidth - 4\tabcolsep) * \real{0.3333}}@{}}
\toprule\noalign{}
\begin{minipage}[b]{\linewidth}\raggedright
Системное требование
\end{minipage} & \begin{minipage}[b]{\linewidth}\raggedright
Программное требование
\end{minipage} & \begin{minipage}[b]{\linewidth}\raggedright
Описательный кейс
\end{minipage} \\
\midrule\noalign{}
\endhead
\bottomrule\noalign{}
\endlastfoot
{\href{../04\_Системные\_требования/DOC-05\_Системные\_требования.pdf\#sr-cfg-001}{SR-CFG-001}} & {\hyperlink{swr-cfg-001}{SWR-CFG-001}} & {\href{../10\_Верификация\_и\_тестирование/DOC-10\_План\_верификации.pdf\#tc-cfg-001}{TC-CFG-001}} \\
{\href{../04\_Системные\_требования/DOC-05\_Системные\_требования.pdf\#sr-cfg-002}{SR-CFG-002}} & {\hyperlink{swr-cfg-002}{SWR-CFG-002}} & {\href{../10\_Верификация\_и\_тестирование/DOC-10\_План\_верификации.pdf\#tc-cfg-002}{TC-CFG-002}} \\
{\href{../04\_Системные\_требования/DOC-05\_Системные\_требования.pdf\#sr-cfg-003}{SR-CFG-003}} & {\hyperlink{swr-cfg-003}{SWR-CFG-003}} & {\href{../10\_Верификация\_и\_тестирование/DOC-10\_План\_верификации.pdf\#tc-cfg-003}{TC-CFG-003}} \\
{\href{../04\_Системные\_требования/DOC-05\_Системные\_требования.pdf\#sr-cfg-004}{SR-CFG-004}} & {\hyperlink{swr-cfg-004}{SWR-CFG-004}} & {\href{../10\_Верификация\_и\_тестирование/DOC-10\_План\_верификации.pdf\#tc-cfg-004}{TC-CFG-004}} \\
{\href{../04\_Системные\_требования/DOC-05\_Системные\_требования.pdf\#sr-cfg-005}{SR-CFG-005}} & {\hyperlink{swr-cfg-005}{SWR-CFG-005}} & {\href{../10\_Верификация\_и\_тестирование/DOC-10\_План\_верификации.pdf\#tc-cfg-005}{TC-CFG-005}} \\
{\href{../04\_Системные\_требования/DOC-05\_Системные\_требования.pdf\#sr-cfg-006}{SR-CFG-006}} & {\hyperlink{swr-cfg-006}{SWR-CFG-006}} & {\href{../10\_Верификация\_и\_тестирование/DOC-10\_План\_верификации.pdf\#tc-cfg-006}{TC-CFG-006}} \\
{\href{../04\_Системные\_требования/DOC-05\_Системные\_требования.pdf\#sr-cfg-007}{SR-CFG-007}} & {\hyperlink{swr-cfg-007}{SWR-CFG-007}} & {\href{../10\_Верификация\_и\_тестирование/DOC-10\_План\_верификации.pdf\#tc-cfg-007}{TC-CFG-007}} \\
{\href{../04\_Системные\_требования/DOC-05\_Системные\_требования.pdf\#sr-cfg-008}{SR-CFG-008}} & {\hyperlink{swr-cfg-008}{SWR-CFG-008}} & {\href{../10\_Верификация\_и\_тестирование/DOC-10\_План\_верификации.pdf\#tc-cfg-008}{TC-CFG-008}} \\
\end{longtable}
}

Для рассмотрения проверяют доступ к чужой группе и чужому мероприятию, выключенный шаблон, повторяющиеся позиции, пустой черновик, равные и обратные границы окна, точные моменты открытия и закрытия, невыбранный и явно запрещённый повтор, редактирование после публикации и неизменность закреплённого варианта после изменения банка заданий.

Проверка пустого мероприятия выполняется с допустимыми группой и обеими границами окна: иначе отказ может относиться к неверной форме запроса. После добавления позиции отдельно рассматриваются отсутствие выбора повтора и явное \texttt{repeat\allowbreak{}Allowed =\allowbreak{} false}. Правила и порядок отказов уже определены в SWR-CFG-008 и DOC-08.

Ведомость преподавателя уточнена в DOC-08: просрочка определяется на момент чтения с учётом своевременных ответов. Условия публикации в TC-CFG-008 уточнены в редакции DOC-10 от 25.11.2026: указаны обе границы окна и последовательность выбора повтора. Это авторские исправления текста по предложениям {\href{../08\_Протоколы\_формальных\_инспекций/DOC-13.5\_Протокол\_ФИ-АРХ-API.pdf\#замечания-к-разбору}{DOC-13.5}}; решение участников о закрытии замечаний ожидается.

\section{Согласованность редакции и условия передачи}\label{5-ux441ux43eux433ux43bux430ux441ux43eux432ux430ux43dux43dux43eux441ux442ux44c-ux440ux435ux434ux430ux43aux446ux438ux438-ux438-ux443ux441ux43bux43eux432ux438ux44f-ux43fux435ux440ux435ux434ux430ux447ux438}

\hypertarget{5-согласованность-редакции-и-условия-передачи}{}

Восемь определений SWR-CFG сохранены в единой главе. Каталог и модерация различают роли и группы; черновик и публикация различают допустимый пустой состав и обязательный выбор повтора; выдача использует все пары опубликованной редакции, закрепляя выбранную пару за личной задачей. В физической модели этим сведениям соответствуют назначения шаблонов, позиции, \texttt{assessment\_\allowbreak{}variant}, \texttt{problem\_\allowbreak{}session} и \texttt{problem}. Повтор не заменяет прошлое прохождение.

Кандидат v1.0 включает подготовленные SR → SWR → TC и проверяемые границы публикации, прав, сроков, вариантов и восстановления. Для редакторской трассировки подготовлен {\href{../05\_Трассировка/DOC-09\_Вклад\_Самбьюи.pdf\#doc-09-вклад-самбьюи-в-матрицу-трассировки}{вклад Самбьюи}}. Кейсы имеют описательный характер: результаты исполнения не заявлены.

До принятия требуется независимая проверка конфигураций назначенным рецензентом, разбор реальных замечаний ФИ-СТ при наличии протокола DOC-13.3 и решение участников. Подготовка кандидата не подтверждает выполнение задачи 3.1: зарегистрированные замечания ФИ-СТ в доступном комплекте отсутствуют. Собственные разделы автор не оценивает как независимый рецензент.

\section{История редакций}\label{ux438ux441ux442ux43eux440ux438ux44f-ux440ux435ux434ux430ux43aux446ux438ux439}

\hypertarget{история-редакций}{}

{\def\LTcaptype{none} % do not increment counter
\begin{longtable}[]{@{}
  >{\raggedright\arraybackslash}p{(\linewidth - 2\tabcolsep) * \real{0.3200}}
  >{\raggedright\arraybackslash}p{(\linewidth - 2\tabcolsep) * \real{0.6800}}@{}}
\toprule\noalign{}
\begin{minipage}[b]{\linewidth}\raggedright
Дата
\end{minipage} & \begin{minipage}[b]{\linewidth}\raggedright
Изменение
\end{minipage} \\
\midrule\noalign{}
\endhead
\bottomrule\noalign{}
\endlastfoot
20.11.2026 & В сводной редакции DOC-06 подготовлены SWR-CFG-001--SWR-CFG-008 \\
25.11.2026 & Самбьюи А.Ж.: подготовлена самостоятельная редакция DOC-06.6 v0.8, общая глава определений, границы подсистемы и таблица SR → SWR → TC. Предметные правила и ID сохранены \\
25.11.2026 & Самбьюи А.Ж.: подготовлен кандидат v1.0 для задачи 3.2; сверены физическая модель, API, кейсы и условия передачи на инспекцию. Восемь SWR-CFG и их ID сохранены \\
\end{longtable}
}

\section{Конфигурации, контрольные мероприятия и модерация}\label{ux43aux43eux43dux444ux438ux433ux443ux440ux430ux446ux438ux438-ux43aux43eux43dux442ux440ux43eux43bux44cux43dux44bux435-ux43cux435ux440ux43eux43fux440ux438ux44fux442ux438ux44f-ux438-ux43cux43eux434ux435ux440ux430ux446ux438ux44f}

\hypertarget{конфигурации-контрольные-мероприятия-и-модерация}{}

Шаблон набора задач --- именованный перечень видов заданий в порядке выдачи. По нему студент может начать самостоятельную практику. Для контрольного мероприятия преподаватель составляет работу, назначает её конкретной группе и публикует в расписании. Позиция --- одно место в составе работы. Вариант --- конкретное условие задания для этой позиции с данными для его проверки. Личная сессия связывает выданные варианты с конкретным студентом и его результатами.

Состояние \texttt{DRAFT} обозначает неопубликованный черновик, доступный преподавателю для редактирования. \texttt{PUBLISHED} обозначает опубликованное мероприятие с закреплённым составом. Факт публикации позволяет группе увидеть назначение в расписании. Возможность начать работу сервер проверяет отдельно. Допуск определяют {\href{../02\_Техническое\_задание/DOC-02\_Техническое\_задание.pdf\#правила-доступа-и-прохождения-контрольного-мероприятия}{правила доступа и прохождения контрольного мероприятия}}.

\subsection{SWR-CFG-001. Выборка шаблонов своей группы}\label{swr-cfg-001-ux432ux44bux431ux43eux440ux43aux430-ux448ux430ux431ux43bux43eux43dux43eux432-ux441ux432ux43eux435ux439-ux433ux440ux443ux43fux43fux44b}

\hypertarget{swr-cfg-001-выборка-шаблонов-своей-группы}{}

\hypertarget{swr-cfg-001}{}

Обработчик каталога должен сверять запрошенную группу с группой подтверждённого пользователя. Пользователю с ролью \texttt{USER} доступны только шаблоны набора задач его группы. Для каждого назначения обработчик возвращает имя шаблона, список типов задач в заданном порядке и признак \texttt{enabled} именно для этой группы. Отключённые шаблоны остаются видны в каталоге группы. Компонент создания сессии принимает только включённый шаблон.

Основание --- {\href{../04\_Системные\_требования/DOC-05\_Системные\_требования.pdf\#sr-cfg-001}{SR-CFG-001}}. Если группе назначены включённый и отключённый шаблоны, каталог различает оба. Создание по включённому шаблону сохраняет число позиций, включая повторяющиеся типы. Запрос каталога чужой группы не раскрывает её шаблоны.

\subsection{SWR-CFG-002. Чтение сессий всех студентов}\label{swr-cfg-002-ux447ux442ux435ux43dux438ux435-ux441ux435ux441ux441ux438ux439-ux432ux441ux435ux445-ux441ux442ux443ux434ux435ux43dux442ux43eux432}

\hypertarget{swr-cfg-002-чтение-сессий-всех-студентов}{}

\hypertarget{swr-cfg-002}{}

Обработчик модерации должен допускать к общему перечню сессий только роль \texttt{ADMIN}. Он читает сессии всех студентов с их владельцами, состояниями сессий и состояниями отдельных задач. Обращение к перечню не запускает предметную проверку и не предоставляет права отправить решение за владельца. Преподавательская роль сама по себе не предоставляет доступа к этому перечню.

Основание --- {\href{../04\_Системные\_требования/DOC-05\_Системные\_требования.pdf\#sr-cfg-002}{SR-CFG-002}}. Если сессии есть у двух студентов из разных групп, в общем перечне модератора видны обе. При обращении только с ролью \texttt{USER} сервер отказывает до выдачи перечня.

\subsection{SWR-CFG-003. Состав контрольного мероприятия}\label{swr-cfg-003-ux441ux43eux441ux442ux430ux432-ux43aux43eux43dux442ux440ux43eux43bux44cux43dux43eux433ux43e-ux43cux435ux440ux43eux43fux440ux438ux44fux442ux438ux44f}

\hypertarget{swr-cfg-003-состав-контрольного-мероприятия}{}

\hypertarget{swr-cfg-003}{}

При создании контрольного мероприятия обработчик должен сохранить черновик \texttt{DRAFT} с уникальным \texttt{assessment\allowbreak{}Id} и составом, заданным преподавателем. Поле \texttt{repeat\allowbreak{}Allowed} хранит выбор преподавателя о повторном прохождении. В черновике его можно отложить. При создании или полной замене черновика пропуск поля и \texttt{null} означают, что выбор ещё не сделан. Значения \texttt{true} и \texttt{false} явно разрешают или запрещают повтор. Другие типы дают HTTP 400. Черновик допускает пустой список позиций. При редактировании своего черновика преподаватель может изменить состав и этот выбор. Для каждой позиции преподаватель выбирает подключённый программный вид либо своё авторское задание с готовыми вариантами. В одной работе допускаются позиции разных видов и несколько позиций одного вида. Преподаватель может просмотреть свои мероприятия и заменить состав своего черновика. Если к моменту записи правки его работа уже опубликована, обработчик отклоняет правку и сообщает автору о публикации. При выдаче личного набора подсистема сессий получает позиции выбранного мероприятия в заданном порядке. Создание черновика не открывает студентам сессии.

Основание --- {\href{../04\_Системные\_требования/DOC-05\_Системные\_требования.pdf\#sr-cfg-003}{SR-CFG-003}}. Для мероприятия из ранцевой задачи и задачи Эль-Гамаля преподаватель видит обе позиции в заданном порядке. При формировании личного набора этим позициям соответствуют две задачи своих типов.

\subsection{SWR-CFG-004. Назначение конкретной группе}\label{swr-cfg-004-ux43dux430ux437ux43dux430ux447ux435ux43dux438ux435-ux43aux43eux43dux43aux440ux435ux442ux43dux43eux439-ux433ux440ux443ux43fux43fux435}

\hypertarget{swr-cfg-004-назначение-конкретной-группе}{}

\hypertarget{swr-cfg-004}{}

При публикации обработчик должен проверить роль \texttt{TEACHER} и действующее назначение указанной группы этому преподавателю. Публикация требует двух величин --- времени открытия \texttt{opensAt} и времени закрытия \texttt{closesAt}, причём \texttt{opens\allowbreak{}At < closes\allowbreak{}At}. Сервер сравнивает их со своими часами. Переход в \texttt{PUBLISHED}, запись группы, границ, значения \texttt{repeat\allowbreak{}Allowed}, порядка позиций и редакций заданий выполняются в одной транзакции. Опубликованные параметры определяют правила всех прохождений этой работы. При отказе черновик сохраняется, в расписании новая запись не появляется. Разные мероприятия одной группы различаются по \texttt{assessment\allowbreak{}Id}.

Основание --- {\href{../04\_Системные\_требования/DOC-05\_Системные\_требования.pdf\#sr-cfg-004}{SR-CFG-004}}. Публикация на группу, которую администратор не закреплял за преподавателем, не создаёт назначения. Равные границы окна также исключают публикацию.

\subsection{SWR-CFG-005. Начало личной сессии по мероприятию}\label{swr-cfg-005-ux43dux430ux447ux430ux43bux43e-ux43bux438ux447ux43dux43eux439-ux441ux435ux441ux441ux438ux438-ux43fux43e-ux43cux435ux440ux43eux43fux440ux438ux44fux442ux438ux44e}

\hypertarget{swr-cfg-005-начало-личной-сессии-по-мероприятию}{}

\hypertarget{swr-cfg-005}{}

По \texttt{assessment\allowbreak{}Id} подсистема должна проверить публикацию и принадлежность работы группе студента, затем определить по серверным часам, открыт ли старт. Окно открывается и закрывается автоматически. До \texttt{opensAt} сервер возвращает \texttt{409 ASSESSMENT\_\allowbreak{}NOT\_\allowbreak{}OPEN}, с момента \texttt{closesAt} --- \texttt{409 ASSESSMENT\_\allowbreak{}CLOSED}. Внутри окна действует ограничение одной сессии \texttt{NEW} на студента. При наличии такой сессии студент получает \texttt{409 ACTIVE\_\allowbreak{}SESSION\_\allowbreak{}EXISTS} и её UUID. Если прежнее прохождение этой работы завершилось, при \texttt{repeat\allowbreak{}Allowed =\allowbreak{} false} сервер возвращает \texttt{409 ASSESSMENT\_\allowbreak{}RETAKE\_\allowbreak{}DISABLED} с UUID прежней сессии в \texttt{details.\allowbreak{}session\allowbreak{}Id}. При \texttt{repeat\allowbreak{}Allowed =\allowbreak{} true} успешный старт создаёт отдельную сессию с новыми UUID сессии и задач для того же мероприятия. Сессии и ответы предыдущих прохождений сохраняются и доступны владельцу по их UUID.

Для начала контроля запрос содержит идентификатор мероприятия. Параметр \texttt{configuration} относится к самостоятельной практике по шаблону своей группы. В запросе должен быть ровно один из этих параметров. Для практики \texttt{assessment\allowbreak{}Id} созданной сессии равен \texttt{null}. Результат такой сессии не может быть зачтён как участие в мероприятии.

Основание --- {\href{../04\_Системные\_требования/DOC-05\_Системные\_требования.pdf\#sr-cfg-005}{SR-CFG-005}}. Ровно в момент открытия студент своей группы может начать работу при соблюдении остальных условий. На границе закрытия новый старт уже невозможен. Любой отказ в создании сохраняет прежние задачи и результаты.

\subsection{SWR-CFG-006. Авторское задание с готовыми ответами}\label{swr-cfg-006-ux430ux432ux442ux43eux440ux441ux43aux43eux435-ux437ux430ux434ux430ux43dux438ux435-ux441-ux433ux43eux442ux43eux432ux44bux43cux438-ux43eux442ux432ux435ux442ux430ux43cux438}

\hypertarget{swr-cfg-006-авторское-задание-с-готовыми-ответами}{}

\hypertarget{swr-cfg-006}{}

Редактор авторских заданий должен позволять преподавателю ввести варианты без программирования. Для каждого варианта он сохраняет пару строк \texttt{statement} и \texttt{expected\allowbreak{}Answer} --- условие и готовый текстовый ответ. Каждый вариант, доступный для выдачи, имеет оба значения. Преподаватель может выбрать сохранённое задание из своего перечня для позиции вида \texttt{CUSTOM\_\allowbreak{}TEXT}. Ответ студента передаётся строкой \texttt{answer}. Подсистема сессий вызывает проверку через тот же интерфейс, что и для программно добавленного вида задания.

Для изменения своего задания преподаватель передаёт название и варианты по \texttt{PUT /\allowbreak{}api/\allowbreak{}custom-\allowbreak{}tasks/\allowbreak{}\{custom\allowbreak{}Task\allowbreak{}Id\}} в той же форме, что и при создании. Компонент сохраняет новую редакцию для будущих публикаций. Мероприятия, уже опубликованные к моменту изменения, продолжают использовать закреплённую редакцию, даже если ни один студент ещё не начал работу.

Проверяющий получает ответ студента и эталон выданного варианта. У каждой из двух строк он удаляет пробелы по краям и затем сравнивает строки целиком. Регистр букв и внутренние пробелы сохраняются. Например, эталон «Красная книга» и ответ с теми же словами принимаются при любом числе пробелов по краям обеих строк. Ответ с маленькой начальной буквой или с двумя пробелами между словами даёт отрицательный вердикт.

Основание --- {\href{../04\_Системные\_требования/DOC-05\_Системные\_требования.pdf\#sr-cfg-006}{SR-CFG-006}}. Для проверки используют два варианта с разными эталонами. Вердикт относится к эталону того варианта, который получил студент.

\subsection{SWR-CFG-007. Выбор и закрепление авторского варианта}\label{swr-cfg-007-ux432ux44bux431ux43eux440-ux438-ux437ux430ux43aux440ux435ux43fux43bux435ux43dux438ux435-ux430ux432ux442ux43eux440ux441ux43aux43eux433ux43e-ux432ux430ux440ux438ux430ux43dux442ux430}

\hypertarget{swr-cfg-007-выбор-и-закрепление-авторского-варианта}{}

\hypertarget{swr-cfg-007}{}

Для каждой позиции с авторским заданием обработчик выдачи должен случайно выбрать одну из пар той редакции задания, которая была закреплена при публикации, и связать выбор с этой позицией в личной сессии. Последующие изменения банка вариантов не меняют ни опубликованное задание, ни уже выданную пару. Открытое условие входит в набор студента. Основание для проверки и эталонный ответ хранятся на сервере. Их нельзя получить в составе условия, через повторное чтение сессии или через личную историю.

Основание --- {\href{../04\_Системные\_требования/DOC-05\_Системные\_требования.pdf\#sr-cfg-007}{SR-CFG-007}}. После перезапуска чтение возвращает прежнее условие, а проверка использует прежний эталон без нового случайного выбора. Повторяющиеся позиции одного вида имеют собственные привязки к вариантам. При случайном выборе один вариант может встретиться повторно.

\subsection{SWR-CFG-008. Публикация в расписании группы}\label{swr-cfg-008-ux43fux443ux431ux43bux438ux43aux430ux446ux438ux44f-ux432-ux440ux430ux441ux43fux438ux441ux430ux43dux438ux438-ux433ux440ux443ux43fux43fux44b}

\hypertarget{swr-cfg-008-публикация-в-расписании-группы}{}

\hypertarget{swr-cfg-008}{}

Компонент расписания должен включить назначенное группе контрольное мероприятие в расписание студентов этой группы после явного действия преподавателя по публикации. Обработчик публикации проверяет сохранённый список позиций в той же транзакции, в которой закрепляет состав и переводит мероприятие в состояние \texttt{PUBLISHED}. Если список пуст, обработчик возвращает отказ с указанием добавить хотя бы одно задание. При корректных параметрах назначения проверка пустого состава выполняется первой и даёт \texttt{409 ASSESSMENT\_\allowbreak{}EMPTY}, даже если режим повторов ещё не выбран. Для непустого состава без выбора \texttt{repeat\allowbreak{}Allowed} сервер возвращает \texttt{409 ASSESSMENT\_\allowbreak{}RETAKE\_\allowbreak{}POLICY\_\allowbreak{}REQUIRED} и просит преподавателя разрешить либо запретить повтор. В обоих случаях сохраняются состояние \texttt{DRAFT} и прежний идентификатор, запись расписания не создаётся. При успехе опубликованное значение \texttt{repeat\allowbreak{}Allowed} всегда равно \texttt{true} или \texttt{false}. С момента публикации студент своей группы видит мероприятие, оба момента окна и разрешение на повторное прохождение, в том числе до открытия. Чтение расписания не создаёт личную сессию и не выдаёт условия или эталоны. Условия доступны студенту в его личном наборе после допуска к началу работы. Поле \texttt{canStart} учитывает открытое окно, наличие активной сессии и возможность повтора. Если у студента были прохождения этого мероприятия, \texttt{sessionId} указывает на его собственную активную сессию мероприятия либо, при её отсутствии, на последнюю завершённую. Остальные сессии доступны владельцу по своим UUID.

Основание --- {\href{../04\_Системные\_требования/DOC-05\_Системные\_требования.pdf\#sr-cfg-008}{SR-CFG-008}}. Проверку выполняют для двух групп. Сразу после публикации назначенная группа видит запись о мероприятии в расписании. В расписании другой группы эта запись отсутствует. Перед открытием из расписания нельзя получить индивидуальные условия или эталонные ответы.


\end{document}
